\subsection{Problems to Solve}
\label{sec:problems-to-solve}

The existing hand-operated expense reporting system at FSP introduces substantial inefficiencies for both employees and the Finance department. Specifically, employees spend excessive time entering detailed receipt information in Excel spreadsheets. Concurrently, the Finance Department undertakes labour-intensive work in verifying outdated spreadsheets.

Moreover, from FSP's perspective, this manual data entry is prone to human error. According to the company's historical data, approximately 30\% of submitted reports require manual correction due to transcription errors or missing information. Consequently, this not only increases the risk of inaccuracy but also results in expenses being delivered to the wrong person. 

Quantitatively, FSP's Barcelona team of 25 employees departs for the UK 2-4 times annually, generating an average of 8 receipts per person per trip. Each employee spends around 15 minutes on a manual expense report, while the Finance Department (composed of 4 people) takes 20 minutes to review. Collectively, this unproductive task accounts for roughly 175 lost hours per year and prolongs reimbursement cycles by about 18 days.

Therefore, the time consumed in this burdensome activity could be considerably reduced, enabling highly skilled employees to focus on activities that deliver higher value to the company. Hence, there is a need to optimize FSP's personnel investment.

